% Results section (4.1-4.6) of the journal paper, following paper/OUTLINE.md.
%
% Every number comes from the final report and validation tables on OneDrive
% (Final_correct_ results/IMD_Mapping_report 1.pdf and validation_combined.docx,
% 2026-10-23), not from the older report/IMD_Mapping_report.tex.
%
% Figures are placeholders (\figfrom) naming the figure of the final report to
% copy in; replace each with \includegraphics once the final figure files are
% in paper/figs/. Cross-references to Sections 2-3 and 5 (sec:data,
% sec:method, sec:validation, sec:discussion, sec:limitations) resolve once
% this file is \input into the full paper.
%
% Standalone build: compiles with pdfLaTeX (Overleaf default).
\documentclass[11pt,a4paper]{article}

\usepackage[margin=2.2cm]{geometry}
\usepackage[T1]{fontenc}
\usepackage[utf8]{inputenc}
\usepackage{amsmath}
\usepackage{booktabs}
\usepackage{multirow}
\usepackage{graphicx}
\usepackage{xcolor}
\usepackage[font=small,labelfont=bf]{caption}
\usepackage[colorlinks,linkcolor=blue!50!black,citecolor=blue!50!black]{hyperref}
\usepackage[capitalise,nameinlink]{cleveref}

\newcommand{\Rsq}{$R^2$}
\newcommand{\pp}{\,pp}
% Placeholder box for a figure still to be copied from the final report.
\newcommand{\figfrom}[1]{\fbox{\parbox[c][3.2cm][c]{0.8\textwidth}{\centering
  \small\itshape Placeholder: #1}}}

\begin{document}

\setcounter{section}{3}
\section{Results}
\label{sec:results}

Two sign conventions are used throughout. Raster-wide and plot-wise
differences are given as \emph{predicted $-$ reference}, so a positive
difference means the map lies above the reference. The bias of the
independent validation follows the definition of
Section~\ref{sec:validation} and is \emph{observed $-$ predicted}, so a
negative bias means the map over-predicts the photo-interpreted
imperviousness. Continuous errors are in percentage points (pp) of
imperviousness on the 0--100 scale. Within each validation table, maps are
ordered by the mean rank of \Rsq{}, Cohen's $\kappa$ and the quadratic
weighted kappa (QWK).

% ---------------------------------------------------------------------------
\subsection{Spatial cross-validation and model robustness}
\label{sec:res-cv}

\Cref{tab:cv} reports the spatially blocked five-fold cross-validation used
to tune every random forest. For each block size it gives the mean RMSE of
the best of the forty hyperparameter settings; the last column gives the RMSE
of the out-of-fold predictions of the selected model. These scores measure
how well each model reproduces its own training target, not the
photo-interpreted ground truth, and serve only for model selection.

\begin{table}[htbp]
  \centering
  \small
  \caption{Random forest tuning by spatially blocked five-fold
  cross-validation (pp). For each block size, the mean RMSE over the five
  rounds of the best of the forty settings; bold marks the selected
  (lowest) block size. Out-of-fold RMSE is that of the selected model.
  Spread is the largest minus the smallest of the three block-size scores.
  Scores are against each run's own training target. Zero-shot maps reuse
  the Milan CLMS-trained models and have no row of their own.}
  \label{tab:cv}
  \begin{tabular}{llccccc}
    \toprule
    & & \multicolumn{3}{c}{Best setting, mean RMSE} & Out-of-fold & \\
    \cmidrule(lr){3-5}
    Run & Predictor & 500\,m & 1{,}000\,m & 2{,}000\,m & RMSE & Spread \\
    \midrule
    \multirow{4}{*}{Milan, CLMS}
      & AlphaEarth    & 13.240 & 13.253 & \textbf{13.210} & 13.24 & 0.04 \\
      & S2 stack      & \textbf{10.503} & 10.504 & 10.530 & 10.51 & 0.03 \\
      & S2 median     & \textbf{10.785} & 10.810 & 10.801 & 10.79 & 0.03 \\
      & S2 percentile & 9.169 & \textbf{9.108} & 9.174 & 9.11 & 0.07 \\
    \midrule
    \multirow{4}{*}{Milan, GHSL}
      & AlphaEarth    & \textbf{18.684} & 18.743 & 18.696 & 18.70 & 0.06 \\
      & S2 stack      & 18.212 & 18.227 & \textbf{18.190} & 18.20 & 0.04 \\
      & S2 median     & 18.744 & \textbf{18.648} & 18.807 & 18.67 & 0.16 \\
      & S2 percentile & 18.180 & \textbf{18.149} & 18.212 & 18.16 & 0.06 \\
    \midrule
    \multirow{3}{*}{Hanoi, local}
      & AlphaEarth    & 22.479 & 22.890 & \textbf{22.448} & 22.54 & 0.44 \\
      & S2 median     & \textbf{23.955} & 24.455 & 23.979 & 24.01 & 0.50 \\
      & S2 stack      & \textbf{24.043} & 24.432 & 24.063 & 24.05 & 0.39 \\
    \midrule
    \multirow{3}{*}{HCMC, local}
      & AlphaEarth    & 21.537 & \textbf{21.512} & 21.646 & 21.55 & 0.13 \\
      & S2 median     & 23.388 & 23.325 & \textbf{23.283} & 23.24 & 0.11 \\
      & S2 stack      & 23.208 & 23.199 & \textbf{23.120} & 23.08 & 0.09 \\
    \bottomrule
  \end{tabular}
\end{table}

\paragraph{Block size has little effect.} For every model, the three block
sizes give almost the same RMSE. The spread is at most 0.16\pp{} in Milan and
at most 0.50\pp{} in Hanoi and HCMC, which is small next to the differences
between predictors and between training targets. No block size is
consistently best: each of 500, 1{,}000 and 2{,}000\,m is selected for some
models. Model selection is therefore insensitive to the spatial scale of the
blocks, and the selected hyperparameters do not depend on that choice.

\paragraph{Cross-validation error depends mostly on the training target.}
The out-of-fold RMSE rises from 9.1--13.2\pp{} for the Milan CLMS-trained
models to 18.2--18.7\pp{} for the Milan GHSL-trained models and
21.6--24.1\pp{} for the locally trained Vietnamese models. In Milan, where
both targets share the same predictors and area, replacing CLMS with GHSL
raises the cross-validation error by 5.5 to 9.1\pp{}: the roofed-surface
target of GHSL is harder to reproduce from 10\,m predictors than the
sealed-surface target of CLMS.

\paragraph{Predictor rankings differ between cross-validation and
independent validation.} Against CLMS, the Sentinel-2 percentile composite
has the lowest cross-validation error (9.11\pp) and AlphaEarth the highest
(13.24\pp). In Vietnam the order reverses: AlphaEarth has the lowest error in
both Hanoi (22.54\pp) and HCMC (21.55\pp), 1.5 to 1.7\pp{} below the
Sentinel-2 composites. As the following sections show, the independent
validation does not reproduce the Milan cross-validation ranking: AlphaEarth
scores above the S2 median against the EarthLabel plots, and the
cross-validation errors (9--24\pp) are well below the independent RMSEs
(26--40\pp). Cross-validation against the training product is thus a guide
for tuning, not a measure of mapping accuracy.

% ---------------------------------------------------------------------------
\subsection{Milan: AlphaEarth versus Sentinel-2}
\label{sec:res-milan}

\begin{figure}[htbp]
  \centering
  \figfrom{final report Fig.~11 --- the eight Milan maps and their
  difference from the benchmark they were trained on}
  \caption{Milan: the eight predicted maps (top two rows trained on CLMS,
  bottom two on GHSL) and their difference maps, predicted map $-$ training
  benchmark. Red: map above the benchmark; blue: below.}
  \label{fig:milan-maps}
\end{figure}

\paragraph{Agreement with the benchmark over the raster.} Every Milan map lies
above its training benchmark on average (\cref{fig:milan-maps},
\cref{tab:milan-diff}). Among the CLMS-trained maps, AlphaEarth has the
smallest mean difference (+4.27\pp) and the S2 percentile composite the
smallest mean absolute difference (7.40\pp) and standard deviation
(10.99\pp). The S2 median departs most from CLMS, with a mean difference of
+8.50\pp{} and a mean absolute difference of 11.21\pp. The stack maps cover
slightly fewer pixels than the others, since the stack has no value where any
of its dates is cloudy.

\begin{table}[htbp]
  \centering
  \small
  \caption{Milan: difference between each predicted map and its training
  benchmark over the whole raster (pp; predicted $-$ benchmark).}
  \label{tab:milan-diff}
  \begin{tabular}{llccc}
    \toprule
    Benchmark & Predicted map & Mean diff. & Mean abs. diff. & SD \\
    \midrule
    \multirow{4}{*}{CLMS}
      & S2 stack      & +5.53 & 8.07 & 11.91 \\
      & S2 percentile & +5.19 & 7.40 & 10.99 \\
      & S2 median     & +8.50 & 11.21 & 13.66 \\
      & AlphaEarth    & +4.27 & 7.82 & 11.28 \\
    \midrule
    \multirow{4}{*}{GHSL}
      & S2 stack      & +6.02 & 8.01 & 12.81 \\
      & S2 percentile & +5.20 & 7.34 & 13.00 \\
      & S2 median     & +6.77 & 9.16 & 13.81 \\
      & AlphaEarth    & +5.47 & 7.41 & 11.88 \\
    \bottomrule
  \end{tabular}
\end{table}

\begin{table}[htbp]
  \centering
  \small
  \setlength{\tabcolsep}{4.5pt}
  \caption{Milan: independent validation against the 450 EarthLabel plots.
  RMSE, MAE and bias in pp; bias is observed $-$ predicted. ``CLMS'' and
  ``GHSL'' alone are the benchmark products scored as maps (grey). Within
  each block, rows run from best to worst by mean rank of \Rsq, $\kappa$ and
  QWK; bold marks the best value in each block.}
  \label{tab:milan-val}
  \begin{tabular}{lccccccccc}
    \toprule
    & \multicolumn{4}{c}{A $\cdot$ continuous error}
    & \multicolumn{2}{c}{B $\cdot$ binary at 50\%}
    & \multicolumn{3}{c}{C $\cdot$ ordinal, 10 levels} \\
    \cmidrule(lr){2-5}\cmidrule(lr){6-7}\cmidrule(lr){8-10}
    Predicted map & RMSE & MAE & Bias & \Rsq & OA & $\kappa$
      & OA exact & OA $\pm$1 & QWK \\
    \midrule
    \multicolumn{10}{l}{\emph{Trained on CLMS}} \\
    S2 stack      & \textbf{25.9} & 17.5 & $-$4.3 & \textbf{0.630} & \textbf{0.840} & \textbf{0.667} & 0.364 & 0.620 & \textbf{0.780} \\
    S2 percentile & 26.1 & 17.8 & $-$4.6 & 0.622 & 0.838 & 0.661 & 0.389 & 0.616 & 0.779 \\
    AlphaEarth    & 27.2 & 19.2 & $-$4.1 & 0.590 & 0.829 & 0.643 & 0.360 & 0.553 & 0.742 \\
    \textcolor{gray}{CLMS (benchmark)} & \textcolor{gray}{27.2} & \textcolor{gray}{\textbf{15.2}} & \textcolor{gray}{\textbf{+0.6}} & \textcolor{gray}{0.589} & \textcolor{gray}{0.829} & \textcolor{gray}{0.641} & \textcolor{gray}{\textbf{0.558}} & \textcolor{gray}{\textbf{0.682}} & \textcolor{gray}{0.775} \\
    S2 median     & 27.6 & 19.8 & $-$6.0 & 0.576 & 0.827 & 0.638 & 0.233 & 0.584 & 0.741 \\
    \midrule
    \multicolumn{10}{l}{\emph{Trained on GHSL}} \\
    AlphaEarth    & \textbf{29.0} & 20.6 & \textbf{+7.7} & \textbf{0.533} & \textbf{0.796} & \textbf{0.542} & 0.367 & 0.542 & \textbf{0.676} \\
    S2 stack      & 29.2 & 20.6 & +8.0 & 0.529 & 0.784 & 0.510 & 0.378 & 0.544 & 0.673 \\
    S2 percentile & 29.2 & \textbf{20.3} & +8.3 & 0.529 & 0.778 & 0.496 & 0.378 & 0.540 & 0.672 \\
    S2 median     & 29.5 & 20.9 & +7.8 & 0.516 & 0.782 & 0.508 & 0.344 & 0.544 & 0.661 \\
    \textcolor{gray}{GHSL (benchmark)} & \textcolor{gray}{37.2} & \textcolor{gray}{23.8} & \textcolor{gray}{+17.6} & \textcolor{gray}{0.234} & \textcolor{gray}{0.707} & \textcolor{gray}{0.293} & \textcolor{gray}{\textbf{0.442}} & \textcolor{gray}{\textbf{0.556}} & \textcolor{gray}{0.493} \\
    \bottomrule
  \end{tabular}
\end{table}

\paragraph{Independent validation.} Against the EarthLabel plots
(\cref{tab:milan-val}), the CLMS-trained Sentinel-2 stack agrees best with
the photo-interpreted imperviousness on all three techniques: RMSE
25.9\pp, \Rsq{} 0.630, $\kappa$ 0.667 and QWK 0.780. The percentile
composite follows within 0.01 on every agreement score. AlphaEarth comes
third, with \Rsq{} 0.590, $\kappa$ 0.643 and QWK 0.742, ahead of the S2
median (\Rsq{} 0.576). Relative to the stack, AlphaEarth is lower by 0.040 in
\Rsq, 0.024 in $\kappa$ and 0.038 in QWK, and its RMSE is 1.3\pp{} higher.

Among the Sentinel-2 composites, accuracy increases with the temporal
information kept per pixel: from the median (10 bands, \Rsq{} 0.576) to the
percentile composite (50 bands, 0.622) and the four-date stack (40 bands,
0.630). The single-value median gains least from a full year of Sentinel-2
dates, whereas the composites that preserve the seasonal spread or individual
dates outperform AlphaEarth, whose 64 bands also summarise a full year.

The stack and the percentile composite agree with the photo-interpreted
plots better than CLMS itself on \Rsq, $\kappa$ and QWK, and have a lower
RMSE. CLMS keeps the lowest MAE (15.2\pp), the smallest bias
(+0.6\pp) and the highest exact and $\pm$1-level ordinal accuracies (0.558
and 0.682). The predicted maps therefore track the ground-truth ranking of
imperviousness as well as or better than CLMS, while CLMS more often hits
the exact level of a plot. All CLMS-trained maps carry a negative bias of
$-$4.1 to $-$6.0\pp, so they slightly over-predict imperviousness, consistent
with their position above CLMS over the raster (\cref{tab:milan-diff}).

When the models are trained on GHSL, the predictor ranking compresses:
the four GHSL-trained maps span only 0.017 in \Rsq{} (0.516--0.533),
against 0.054 for the CLMS-trained maps, and AlphaEarth moves to first place
on \Rsq, $\kappa$ and QWK. With the weaker training target, the choice of
predictor makes little difference.

% ---------------------------------------------------------------------------
\subsection{Effect of the training target: CLMS versus GHS-BUILT-S}
\label{sec:res-target}

Milan is the only city where both training targets are available, which
allows the effect of the target to be isolated with the predictors, sample
design and validation plots held fixed. This mimics the Vietnamese setting,
where CLMS does not exist and GHSL is the only option.

\paragraph{The target matters more than the predictor.} Every CLMS-trained
map scores higher than every GHSL-trained map on \Rsq, $\kappa$ and QWK,
whatever the predictor (\cref{tab:milan-val}). Switching the target from
CLMS to GHSL lowers \Rsq{} by 0.101 for the stack, 0.093 for the percentile
composite, 0.060 for the median and 0.057 for AlphaEarth. These losses are
larger than the largest difference between predictors trained on the same
target (0.054). The bias changes sign: from $-$4.1 to $-$6.0\pp{}
(over-prediction) with CLMS to +7.7 to +8.3\pp{} (under-prediction) with
GHSL.

\paragraph{Raster-wide agreement hides a spatial shift.} Over the whole
raster, the GHSL-trained maps are close to CLMS on average, with mean
differences of $-$1.57 to $-$3.14\pp{} (\cref{tab:ghsl-clms}). The
difference maps (\cref{fig:ghsl-clms}) show that this near-zero mean
results from opposite errors that cancel: the GHSL-trained maps lie below
CLMS in the urban core and above it in rural areas. Both patterns follow
from what GHSL measures. It counts roofed surface only, so roads, car parks
and other sealed open surfaces, which are dense in the city, are learned as
pervious.

\begin{figure}[htbp]
  \centering
  \figfrom{final report Fig.~12 --- the four GHSL-trained maps and CLMS,
  and their difference from CLMS over the raster}
  \caption{Milan: the four GHSL-trained maps and CLMS (top), and their
  difference maps, GHSL-trained map $-$ CLMS (bottom).}
  \label{fig:ghsl-clms}
\end{figure}

\begin{table}[htbp]
  \centering
  \small
  \caption{Milan: GHSL-trained maps $-$ CLMS over the whole raster (pp).}
  \label{tab:ghsl-clms}
  \begin{tabular}{lccc}
    \toprule
    Predicted map & Mean diff. & Mean abs. diff. & SD \\
    \midrule
    S2 stack $\cdot$ GHSL      & $-$1.97 & 9.10 & 14.52 \\
    S2 percentile $\cdot$ GHSL & $-$3.14 & 8.23 & 13.65 \\
    S2 median $\cdot$ GHSL     & $-$1.57 & 10.15 & 15.38 \\
    AlphaEarth $\cdot$ GHSL    & $-$2.86 & 8.59 & 14.34 \\
    \bottomrule
  \end{tabular}
\end{table}

\paragraph{At the validation plots.} The 450 plots are stratified across
imperviousness classes, so intermediate and high imperviousness carry more
weight than over the raster, and the shift becomes clear
(\cref{tab:milan-plots}, \cref{fig:milan-plots}). All four GHSL-trained
maps lie below CLMS by $-$7.10 to $-$7.75\pp{} and below the photo-interpreted
values by $-$7.67 to $-$8.32\pp. The CLMS-trained maps, by contrast, lie
above the photo-interpreted values by +4.06 to +6.01\pp.

\begin{figure}[htbp]
  \centering
  \figfrom{final report Fig.~14 --- CLMS-trained and GHSL-trained maps
  against the 450 photo-interpreted plots}
  \caption{Milan: CLMS-trained (top two rows) and GHSL-trained (bottom two
  rows) maps against the 450 photo-interpreted plots: each map with the plot
  locations, and the difference, map $-$ photo-interpreted value, at each
  plot.}
  \label{fig:milan-plots}
\end{figure}

\begin{table}[htbp]
  \centering
  \small
  \caption{Milan: difference statistics at the 450 plots (pp; predicted
  $-$ reference, where the reference is CLMS or the photo-interpreted value).
  The sign is opposite to that of the bias in \cref{tab:milan-val}.}
  \label{tab:milan-plots}
  \begin{tabular}{llcccc}
    \toprule
    Comparison & Predictor & Mean diff. & Mean abs. diff. & SD & RMSE \\
    \midrule
    \multirow{4}{*}{GHSL-trained vs CLMS}
      & S2 stack      & $-$7.43 & 15.44 & 19.35 & 20.71 \\
      & S2 percentile & $-$7.75 & 15.11 & 19.08 & 20.57 \\
      & S2 median     & $-$7.27 & 15.98 & 19.81 & 21.08 \\
      & AlphaEarth    & $-$7.10 & 14.64 & 18.74 & 20.03 \\
    \midrule
    \multirow{4}{*}{CLMS-trained vs plots}
      & S2 stack      & +4.31 & 17.48 & 25.52 & 25.85 \\
      & S2 percentile & +4.58 & 17.76 & 25.75 & 26.13 \\
      & S2 median     & +6.01 & 19.82 & 27.02 & 27.65 \\
      & AlphaEarth    & +4.06 & 19.18 & 26.94 & 27.21 \\
    \midrule
    \multirow{4}{*}{GHSL-trained vs plots}
      & S2 stack      & $-$8.00 & 20.57 & 28.08 & 29.16 \\
      & S2 percentile & $-$8.32 & 20.26 & 27.97 & 29.15 \\
      & S2 median     & $-$7.83 & 20.89 & 28.52 & 29.55 \\
      & AlphaEarth    & $-$7.67 & 20.63 & 28.04 & 29.04 \\
    \bottomrule
  \end{tabular}
\end{table}

\paragraph{Cost of replacing CLMS with GHSL.} \Cref{tab:target-change}
quantifies the change at the plots. Replacing CLMS with GHSL as the training
target raises the RMSE by 1.83 to 3.31\pp{} and moves the mean difference by
$-$11.73 to $-$13.84\pp, from over-prediction to under-prediction. The RMSE
increase is largest for the stack (+3.31\pp) and the percentile composite
(+3.02\pp), the two best CLMS-trained maps, and smallest for AlphaEarth
(+1.83\pp). The predictors that exploit a good target best lose most when
the target is degraded.

\begin{table}[htbp]
  \centering
  \small
  \caption{Milan: change when GHSL replaces CLMS as the training target,
  against the 450 photo-interpreted plots (pp): GHSL-trained row minus
  CLMS-trained row of \cref{tab:milan-plots}.}
  \label{tab:target-change}
  \begin{tabular}{lcccc}
    \toprule
    Predictor & $\Delta$ Mean diff. & $\Delta$ Mean abs. diff. & $\Delta$ SD
      & $\Delta$ RMSE \\
    \midrule
    S2 stack      & $-$12.31 & +3.09 & +2.56 & +3.31 \\
    S2 percentile & $-$12.90 & +2.50 & +2.22 & +3.02 \\
    S2 median     & $-$13.84 & +1.07 & +1.50 & +1.90 \\
    AlphaEarth    & $-$11.73 & +1.45 & +1.10 & +1.83 \\
    \bottomrule
  \end{tabular}
\end{table}

\paragraph{GHSL-trained models improve on GHSL itself.} Although GHSL is a
weaker target, the models trained on it agree with the photo-interpreted
plots far better than the GHSL product does. GHSL scores lowest of all ten
Milan maps (\Rsq{} 0.234, $\kappa$ 0.293, QWK 0.493) with a bias of
+17.6\pp, whereas the GHSL-trained models reach \Rsq{} 0.516--0.533 and
roughly halve the bias, to +7.7--8.3\pp. A regression model fitted to GHSL
thus recovers part of the imperviousness that GHSL leaves out by design.
This is the situation of the locally retrained Vietnamese models, which
are trained on GHSL alone.

% ---------------------------------------------------------------------------
\subsection{Hanoi: local retraining versus zero-shot transfer}
\label{sec:res-hanoi}

In Hanoi, five maps are compared (\cref{fig:hanoi-maps}): three locally
retrained maps (AlphaEarth, S2 median and S2 stack, trained on GHSL points
from Hanoi) and two zero-shot maps (AlphaEarth and S2 median, from the Milan
CLMS-trained models applied unchanged). The Hanoi stack combines three
dates (30 bands).

\begin{figure}[htbp]
  \centering
  \figfrom{final report Fig.~15 --- the five Hanoi maps and GHSL, and their
  difference from GHSL}
  \caption{Hanoi: the five predicted maps and GHSL (top), and their
  difference maps, predicted map $-$ GHSL (bottom). Red: map above GHSL.}
  \label{fig:hanoi-maps}
\end{figure}

\begin{table}[htbp]
  \centering
  \small
  \caption{Hanoi and HCMC: difference between each predicted map and GHSL
  over the whole raster (pp; predicted $-$ GHSL).}
  \label{tab:viet-diff}
  \begin{tabular}{llccc}
    \toprule
    City & Predicted map & Mean diff. & Mean abs. diff. & SD \\
    \midrule
    \multirow{5}{*}{Hanoi}
      & AlphaEarth, local retrain & +7.64 & 11.23 & 15.37 \\
      & AlphaEarth, zero-shot     & +33.65 & 34.08 & 18.06 \\
      & S2 median, local retrain  & +10.31 & 14.52 & 17.95 \\
      & S2 median, zero-shot      & +36.66 & 37.65 & 25.89 \\
      & S2 stack, local retrain   & +10.37 & 14.47 & 17.24 \\
    \midrule
    \multirow{5}{*}{HCMC}
      & AlphaEarth, local retrain & +7.64 & 13.66 & 17.54 \\
      & AlphaEarth, zero-shot     & +45.00 & 45.36 & 20.55 \\
      & S2 median, local retrain  & +9.18 & 16.36 & 20.08 \\
      & S2 median, zero-shot      & +29.39 & 31.13 & 26.46 \\
      & S2 stack, local retrain   & +9.17 & 16.22 & 19.88 \\
    \bottomrule
  \end{tabular}
\end{table}

\paragraph{Agreement with GHSL over the raster.} Both zero-shot maps lie far
above GHSL, with mean differences of +33.65\pp{} for AlphaEarth and
+36.66\pp{} for the S2 median (\cref{tab:viet-diff}). Local retraining
reduces this to +7.64\pp{} for AlphaEarth, +10.31\pp{} for the S2 median and
+10.37\pp{} for the S2 stack. Part of the zero-shot excess is expected, since
the Milan models were trained on CLMS, which includes roads and pavements
that GHSL omits; but its size, between three and more than four times
that of the locally retrained maps, points to a level shift between the cities.

\begin{table}[htbp]
  \centering
  \small
  \setlength{\tabcolsep}{4.5pt}
  \caption{Hanoi: independent validation against the 450 EarthLabel plots.
  Units, sign and row order as in \cref{tab:milan-val}.}
  \label{tab:hanoi-val}
  \begin{tabular}{lccccccccc}
    \toprule
    & \multicolumn{4}{c}{A $\cdot$ continuous error}
    & \multicolumn{2}{c}{B $\cdot$ binary at 50\%}
    & \multicolumn{3}{c}{C $\cdot$ ordinal, 10 levels} \\
    \cmidrule(lr){2-5}\cmidrule(lr){6-7}\cmidrule(lr){8-10}
    Predicted map & RMSE & MAE & Bias & \Rsq & OA & $\kappa$
      & OA exact & OA $\pm$1 & QWK \\
    \midrule
    AlphaEarth, local retrain & \textbf{28.4} & \textbf{21.8} & \textbf{+6.0} & \textbf{0.617} & 0.873 & 0.742 & 0.296 & 0.449 & \textbf{0.731} \\
    S2 stack, local retrain   & 29.3 & 23.2 & +6.8 & 0.591 & \textbf{0.880} & \textbf{0.755} & 0.202 & 0.444 & 0.695 \\
    S2 median, local retrain  & 30.0 & 23.5 & +7.3 & 0.571 & 0.856 & 0.706 & 0.216 & 0.438 & 0.690 \\
    AlphaEarth, zero-shot     & 34.8 & 28.5 & $-$16.2 & 0.425 & 0.827 & 0.660 & 0.033 & 0.338 & 0.598 \\
    S2 median, zero-shot      & 35.0 & 26.0 & $-$16.6 & 0.416 & 0.789 & 0.587 & 0.198 & 0.447 & 0.617 \\
    \textcolor{gray}{GHSL (benchmark)} & \textcolor{gray}{40.0} & \textcolor{gray}{26.7} & \textcolor{gray}{+18.0} & \textcolor{gray}{0.239} & \textcolor{gray}{0.736} & \textcolor{gray}{0.444} & \textcolor{gray}{\textbf{0.376}} & \textcolor{gray}{\textbf{0.527}} & \textcolor{gray}{0.513} \\
    \bottomrule
  \end{tabular}
\end{table}

\paragraph{Independent validation.} Local retraining outperforms zero-shot
transfer on all three techniques (\cref{tab:hanoi-val}): each of the three
locally retrained maps scores higher than both zero-shot maps on \Rsq,
$\kappa$ and QWK. AlphaEarth with local retraining is the best Hanoi map on
the continuous and ordinal techniques, with RMSE 28.4\pp, \Rsq{} 0.617 and
QWK 0.731. The locally retrained stack has the highest binary agreement
($\kappa$ 0.755, OA 0.880), only 0.013 in $\kappa$ above AlphaEarth. Among
the local models, AlphaEarth exceeds the S2 median by 0.046 in \Rsq{} and
0.041 in QWK.

For AlphaEarth, moving from zero-shot transfer to local retraining raises
\Rsq{} from 0.425 to 0.617 and QWK from 0.598 to 0.731, and lowers the RMSE
by 6.4\pp. Both zero-shot maps over-predict strongly, with biases of
$-$16.2 and $-$16.6\pp. The AlphaEarth zero-shot map in particular places
almost no plot on its correct ordinal level (exact OA 0.033), even though
its binary accuracy remains high (0.827): it separates impervious from
pervious plots at the 50\% threshold but at a systematically inflated
level. The locally retrained maps instead under-predict, by +6.0 to
+7.3\pp, inheriting part of the roof-only bias of GHSL, which itself
under-predicts by +18.0\pp{} and scores lowest of the six maps on \Rsq{}
(0.239) and $\kappa$ (0.444).

% ---------------------------------------------------------------------------
\subsection{Ho Chi Minh City: local retraining versus zero-shot transfer}
\label{sec:res-hcmc}

The same five maps are produced for HCMC (\cref{fig:hcmc-maps}). Here only
one Sentinel-2 date passed the screening for the stack, which is therefore a
single clear image of 10 bands.

\begin{figure}[htbp]
  \centering
  \figfrom{final report Fig.~16 --- the five HCMC maps and GHSL, and their
  difference from GHSL}
  \caption{Ho Chi Minh City: the five predicted maps and GHSL (top), and
  their difference maps, predicted map $-$ GHSL (bottom). Red: map above
  GHSL.}
  \label{fig:hcmc-maps}
\end{figure}

\paragraph{Agreement with GHSL over the raster.} The raster comparison
repeats the Hanoi pattern (\cref{tab:viet-diff}). The zero-shot maps lie far
above GHSL, by +45.00\pp{} for AlphaEarth, the largest departure in the
study, and +29.39\pp{} for the S2 median. Local retraining reduces the mean
difference to +7.64\pp{} for AlphaEarth, +9.18\pp{} for the S2 median and
+9.17\pp{} for the S2 stack.

\begin{table}[htbp]
  \centering
  \small
  \setlength{\tabcolsep}{4.5pt}
  \caption{Ho Chi Minh City: independent validation against the 450
  EarthLabel plots. Units, sign and row order as in \cref{tab:milan-val}.}
  \label{tab:hcmc-val}
  \begin{tabular}{lccccccccc}
    \toprule
    & \multicolumn{4}{c}{A $\cdot$ continuous error}
    & \multicolumn{2}{c}{B $\cdot$ binary at 50\%}
    & \multicolumn{3}{c}{C $\cdot$ ordinal, 10 levels} \\
    \cmidrule(lr){2-5}\cmidrule(lr){6-7}\cmidrule(lr){8-10}
    Predicted map & RMSE & MAE & Bias & \Rsq & OA & $\kappa$
      & OA exact & OA $\pm$1 & QWK \\
    \midrule
    S2 median, zero-shot      & \textbf{27.4} & \textbf{18.0} & \textbf{$-$7.6} & \textbf{0.643} & \textbf{0.847} & \textbf{0.691} & \textbf{0.378} & \textbf{0.633} & \textbf{0.782} \\
    AlphaEarth, local retrain & 28.8 & 22.3 & +11.1 & 0.605 & 0.833 & 0.669 & 0.264 & 0.422 & 0.731 \\
    S2 stack, local retrain   & 29.9 & 23.7 & +10.8 & 0.574 & 0.813 & 0.630 & 0.200 & 0.422 & 0.701 \\
    S2 median, local retrain  & 30.8 & 24.4 & +11.4 & 0.546 & 0.793 & 0.591 & 0.224 & 0.438 & 0.687 \\
    \textcolor{gray}{GHSL (benchmark)} & \textcolor{gray}{37.6} & \textcolor{gray}{26.0} & \textcolor{gray}{+19.0} & \textcolor{gray}{0.327} & \textcolor{gray}{0.729} & \textcolor{gray}{0.466} & \textcolor{gray}{0.324} & \textcolor{gray}{0.491} & \textcolor{gray}{0.593} \\
    AlphaEarth, zero-shot     & 40.3 & 33.0 & $-$21.7 & 0.224 & 0.684 & 0.351 & 0.044 & 0.364 & 0.435 \\
    \bottomrule
  \end{tabular}
\end{table}

\paragraph{Independent validation.} Among the locally retrained maps,
AlphaEarth is again the best, on all three techniques (\cref{tab:hcmc-val}):
\Rsq{} 0.605, $\kappa$ 0.669 and QWK 0.731, ahead of the single-date stack
(\Rsq{} 0.574) and the S2 median (0.546). Its margin over the S2 median,
0.059 in \Rsq{} and 0.044 in QWK, is similar to that in Hanoi. The locally
retrained maps under-predict by +10.8 to +11.4\pp, more than in Hanoi.

The two zero-shot maps diverge sharply. AlphaEarth zero-shot is the weakest
map in HCMC, below GHSL itself, with \Rsq{} 0.224, $\kappa$ 0.351, QWK 0.435
and a bias of $-$21.7\pp. The S2 median zero-shot map, by contrast, is the
best map overall, with RMSE 27.4\pp, \Rsq{} 0.643, $\kappa$ 0.691 and QWK
0.782, and it leads on every metric, including the smallest bias.

\paragraph{The HCMC zero-shot result reflects level matching, not
transfer.} The S2 median zero-shot map carries the imperviousness level it
learned in Milan into both Vietnamese cities (\cref{fig:level}). At the
validation plots its mean predicted imperviousness is 61.0\% in Hanoi and
57.9\% in HCMC, while the mean photo-interpreted value is 44.4\% in Hanoi
and 50.2\% in HCMC. The carried-over level therefore lands 16.61\pp{} above
the reference in Hanoi but only 7.64\pp{} above it in HCMC. The same model
ranks first in HCMC and last of the five Hanoi maps on \Rsq. If its HCMC
score reflected transferable mapping skill, it would also lead in Hanoi; it
reflects instead a city-level offset that happens to match the HCMC
reference.

\begin{figure}[htbp]
  \centering
  \figfrom{final report Fig.~17 --- mean predicted IMD of the five
  Vietnamese maps against the mean photo-interpreted value}
  \caption{Mean predicted imperviousness of the five Vietnamese maps against
  the mean photo-interpreted value (vertical line), in Hanoi (top) and HCMC
  (bottom). Filled markers: mean at the 450 validation plots; open markers:
  mean over the whole raster. The two populations differ because the plots
  are stratified across imperviousness classes.}
  \label{fig:level}
\end{figure}

% ---------------------------------------------------------------------------
\subsection{Cross-city comparison}
\label{sec:res-cross}

\Cref{tab:cross} brings together the main maps of the three cities. Four
patterns hold across them.

\begin{table}[htbp]
  \centering
  \small
  \caption{Cross-city summary of the independent validation (450 EarthLabel
  plots per city). Milan rows are trained on CLMS unless marked GHSL; Hanoi
  and HCMC local models are trained on GHSL. RMSE and bias in pp; bias is
  observed $-$ predicted. Bold marks the best map in each city.}
  \label{tab:cross}
  \begin{tabular}{llccccc}
    \toprule
    City & Map & RMSE & Bias & \Rsq & $\kappa$ & QWK \\
    \midrule
    \multirow{5}{*}{Milan}
      & S2 stack                  & \textbf{25.9} & $-$4.3 & \textbf{0.630} & \textbf{0.667} & \textbf{0.780} \\
      & AlphaEarth                & 27.2 & $-$4.1 & 0.590 & 0.643 & 0.742 \\
      & AlphaEarth $\cdot$ GHSL   & 29.0 & +7.7 & 0.533 & 0.542 & 0.676 \\
      & S2 median $\cdot$ GHSL    & 29.5 & +7.8 & 0.516 & 0.508 & 0.661 \\
      & GHSL (benchmark)          & 37.2 & +17.6 & 0.234 & 0.293 & 0.493 \\
    \midrule
    \multirow{5}{*}{Hanoi}
      & AlphaEarth, local         & \textbf{28.4} & +6.0 & \textbf{0.617} & 0.742 & \textbf{0.731} \\
      & S2 median, local          & 30.0 & +7.3 & 0.571 & 0.706 & 0.690 \\
      & AlphaEarth, zero-shot     & 34.8 & $-$16.2 & 0.425 & 0.660 & 0.598 \\
      & S2 median, zero-shot      & 35.0 & $-$16.6 & 0.416 & 0.587 & 0.617 \\
      & GHSL (benchmark)          & 40.0 & +18.0 & 0.239 & 0.444 & 0.513 \\
    \midrule
    \multirow{5}{*}{HCMC}
      & AlphaEarth, local         & 28.8 & +11.1 & 0.605 & 0.669 & 0.731 \\
      & S2 median, local          & 30.8 & +11.4 & 0.546 & 0.591 & 0.687 \\
      & AlphaEarth, zero-shot     & 40.3 & $-$21.7 & 0.224 & 0.351 & 0.435 \\
      & S2 median, zero-shot      & \textbf{27.4} & $-$7.6 & \textbf{0.643} & \textbf{0.691} & \textbf{0.782} \\
      & GHSL (benchmark)          & 37.6 & +19.0 & 0.327 & 0.466 & 0.593 \\
    \bottomrule
  \end{tabular}
\end{table}

\paragraph{AlphaEarth gives the most consistent locally trained maps.}
Trained locally, AlphaEarth reaches \Rsq{} 0.590 in Milan (CLMS), 0.617 in
Hanoi and 0.605 in HCMC, a range of 0.027, and QWK 0.742, 0.731 and 0.731. It
is the best locally trained map in both Vietnamese cities and the best
GHSL-trained map in Milan. Sentinel-2 depends more on the available
imagery: with 30 clear Milan dates, the stack and percentile composites
outperform AlphaEarth, but with only three or four usable dates in Vietnam
(one for the HCMC stack), every locally trained Sentinel-2 map trails
AlphaEarth, by 0.026
to 0.059 in \Rsq. The advantage of the embeddings is largest where
Sentinel-2 compositing is most constrained by cloud cover.

\paragraph{Local retraining is needed; zero-shot transfer is unreliable.}
AlphaEarth zero-shot transfer fails in both Vietnamese cities (\Rsq{} 0.425
and 0.224), with the largest over-prediction of any map (bias $-$16.2 and
$-$21.7\pp). Sentinel-2 median zero-shot transfer fails in Hanoi (\Rsq{}
0.416) and succeeds in HCMC (0.643) only through level matching
(\cref{sec:res-hcmc}). Zero-shot transfer from Milan thus gives no map that
is reliably usable across both cities, whereas local retraining gives
\Rsq{} 0.55--0.62 in both.

\paragraph{The roof-only target leaves a consistent under-prediction.} The
GHSL product under-predicts the photo-interpreted imperviousness by a
similar amount in all three cities, +17.6, +18.0 and +19.0\pp, and is among
the weakest maps everywhere (\Rsq{} 0.234--0.327). Every model trained on
GHSL, in Milan or locally in Vietnam, improves substantially on GHSL itself
(\Rsq{} 0.516--0.617) and roughly halves its bias, to +6.0 to +11.4\pp, but
none removes it. The Milan comparison (\cref{sec:res-target}) shows that
this residual under-prediction comes from the target, not from the
predictors: training on CLMS instead turns it into a small over-prediction.

\paragraph{Overall accuracy is similar across cities.} The best map of each
city reaches an RMSE of 25.9--28.4\pp{} and \Rsq{} of 0.617--0.643, or
0.605 in HCMC when the level-matched zero-shot map is set aside. With
comparable predictors and a GHSL target, the Vietnamese local models (\Rsq{}
0.617 and 0.605 for AlphaEarth) score slightly higher than their Milan
GHSL-trained counterpart (0.533), although differences in the distribution
of plot values between cities limit direct comparison of absolute scores.

\end{document}
